% Part: proof-theory
% Chapter: sequent-calculus
% Section: interpretation-rules

\documentclass[../../../include/open-logic-section]{subfiles}

\begin{document}

\olfileid[hi]{pt}{seq}{irl}

\olsection{नियमों की व्याख्या}

अनुक्रम की व्याख्या याद करें: अनुक्रम $\Gamma \Sequent \Delta$ किसी
!!a{structure}~$\Struct{M}$ में तब सत्य है जब या तो कम-से-कम एक
$!A \in \Gamma$ असत्य हो, या कम-से-कम एक~$!A \in \Delta$ सत्य हो। हम
\Log{G3c} के तार्किक नियमों को उनके प्रमुख !!{formula}s की सत्यता-शर्तों का
संहिताकरण मान सकते हैं। \RightR{\lif} पर विचार करें। इसका प्रमुख !!{formula}
$!A \lif !B$ है। यह तब सत्य है जब या तो $!A$ असत्य हो अथवा~$!B$ सत्य हो।
अतः अनुक्रम $\ \Sequent !A \lif !B$ सत्य है यदि और केवल यदि $!A \Sequent
!B$ सत्य है। इन अनुक्रमों के पूर्वपक्षों और अनुवर्तियों में प्रसंग के
!!{formula}s $\Gamma$, $\Delta$ जोड़ने पर ठीक \RightR{\lif} नियम का निष्कर्ष
और आधार-वाक्य मिलते हैं। ~\LeftR{\lif} की स्थिति समान है। $!A \lif !B$ तब
असत्य है जब $!A$ सत्य और $!B$ असत्य हो। अतः $!A \lif !B \Sequent \ $ सत्य
है यदि और केवल यदि $\ \Sequent !A$ और $!B \Sequent \ $ दोनों सत्य हों।
\LeftR{\lif} का निष्कर्ष उसके आधार-वाक्यों के संयोजन के समतुल्य है। यह ढाँचा
सुनिश्चित करता है कि \Log{G3c} के नियम निर्दोष हैं (यदि सभी आधार-वाक्य सत्य
हैं तो निष्कर्ष सत्य है)। वे पूर्णता भी सुनिश्चित करते हैं।

वास्तव में, किसी भी सत्यता-फलनीय संयोजक की सत्यता-शर्तें इस प्रकार अनुक्रमों
के समुच्चय के रूप में बताई जा सकती हैं। यह इस तथ्य से मिलता है कि चिरसम्मत
तर्क के प्रत्येक सत्यता-फलन को संयोजक प्रसामान्य रूप के !!a{formula} से बताया
जा सकता है: परमाण्विक और निषेधित परमाण्विक !!{formula}s के वियोजनों का संयोजन।
उदाहरणतः $!A \oplus !B$ को तब और केवल तब सत्य मानें जब $!A$ और $!B$ में से
ठीक एक सत्य हो, अर्थात ``अपवर्जी अथवा''। तब $!A \oplus !B$ सत्य है यदि और
केवल यदि $(!A \lor !B) \land (\lnot !A \lor \lnot !B)$ सत्य है, और वह असत्य
है यदि $(\lnot !A \lor !B) \land (!A \lor \lnot !B)$ सत्य है। अतः हम
$\oplus$ के लिए ये अनुक्रम-कलन नियम प्रस्तुत कर सकते हैं:
\begin{align*}
&
\Axiom$\Gamma \fCenter \Delta, !A, !B$
\Axiom$!A, !B, \Gamma \fCenter \Delta$
\RightLabel{\RightR{\oplus}}
\BinaryInf$\Gamma \fCenter \Delta, !A \oplus !B$
\DisplayProof
\\
& \Axiom$!A, \Gamma \fCenter \Delta, !B$
\Axiom$!B, \Gamma \fCenter \Delta, !A$
\RightLabel{\RightR{\oplus}}
\BinaryInf$!A \oplus !B, \Gamma \fCenter \Delta$
\DisplayProof
\end{align*}
और ये नियम भी निर्दोष और पूर्ण होंगे। इनके निष्कर्ष तब और केवल तब सत्य हैं
जब सभी आधार-वाक्य सत्य हों।

\begin{prob}
मान लें कि $!A \downarrow !B$ तब और केवल तब सत्य है जब न $!A$ सत्य हो और
न~$!B$ (``NOR'')। इसके \Log{G3c} नियम क्या होंगे? (ऐसे दो अनुक्रम खोजें जो
एक साथ तब और केवल तब सत्य हों जब $!A \downarrow !B$ सत्य हो। इससे आपको
\RightR{\downarrow} नियम मिलेगा। फिर ऐसा एक अनुक्रम खोजें जो तब और केवल तब
सत्य हो जब $!A \downarrow !B$ असत्य हो। इससे आपको \LeftR{\downarrow} नियम
मिलेगा।)
\end{prob}

\Log{G3c} में प्रत्येक संयोजक और परिमाणक के ठीक दो नियम हैं: एक वाम नियम और
एक दक्षिण नियम। किंतु \Log{G1c} में दो \LeftR{\land} और दो
\RightR{\lor} नियम हैं। इसका कारण यह है कि जेन्ट्सेन की मूल प्रणाली
\Log{LK}, जिसके आधार पर \Log{G1c} बना है, का एक महत्त्वपूर्ण अचिरसम्मत तर्क,
अर्थात अंतःप्रज्ञात्मक तर्क, के लिए \Log{LJ} नामक रूपांतर था। अंतःप्रज्ञात्मक
तर्क के लिए निर्दोष और पूर्ण प्रणाली पाने के लिए \Log{LJ} के हर अनुक्रम के
अनुवर्ती में अधिकतम एक !!{formula} आने का प्रतिबंध है। इससे \Log{G3c} का
\RightR{\lor} नियम उपयोग करना असंभव हो जाता है, क्योंकि उसके आधार-वाक्य के
अनुवर्ती में $!A$ और~$!B$ दोनों हैं। इसलिए जेन्ट्सेन ने \Log{LJ} के लिए
\RightR{\lor} नियमों का एक युग्म इस्तेमाल किया और \Log{LK} के लिए भी वही
रखा। इसी प्रकार अंतःप्रज्ञात्मक रूपांतर~\Log{G1i}, बस ऐसा~\Log{G1c} है जिसमें
सभी अनुवर्तियों में अधिकतम एक !!{formula} आने का प्रतिबंध है। (\Log{G3c} का
भी एक अंतःप्रज्ञात्मक रूपांतर है, किंतु उसके कुछ नियम \Log{G3c} के अनुरूपी
नियमों से भिन्न हैं। उदाहरणतः उसमें भी \RightR{\lor} नियमों का युग्म है।)

यह देखना सरल है कि \Log{G1c} के संरचनात्मक नियम निर्दोष हैं। उदाहरणतः, यदि
$\Gamma \Sequent \Delta$ में दाईं ओर कोई सत्य !!{formula} अथवा बाईं ओर कोई
असत्य !!{formula} है, तो $\Gamma \Sequent \Delta, !A$ में भी वही है; $!A$
सत्य है या असत्य, इससे कोई अंतर नहीं पड़ता। अतः यदि ~\RightR{\Weakening} का
आधार-वाक्य सत्य है तो उसका निष्कर्ष भी सत्य है। ध्यान दें कि यहाँ इसका विलोम
सत्य नहीं है। निष्कर्ष $!A$ के सत्य होने के कारण सत्य हो सकता है, और तब
आधार-वाक्य के सत्य होने की प्रत्याभूति नहीं है।

संरचनात्मक नियम आखिर क्यों रखें? उदाहरणतः $\ \Sequent !A \lor \lnot !A$
को !!{prove} करने के लिए संकुचन (\RightR{\Contraction},
\LeftR{\Contraction}) चाहिए। यदि आप $!A \Sequent !A$ से आरंभ करें, तो पहले
~\RightR{\lnot} द्वारा $\Sequent !A, \lnot !A$ मिलता है। फिर आप
\RightR{\lor} के दोनों रूप क्रमशः उपयोग कर सकते हैं: एक बार $\lnot !A$ से
$!A \lor \lnot !A$ पाने के लिए, और एक बार~$!A$ से $!A \lor \lnot !A$ पाने
के लिए। इससे $\ \Sequent !A \lor \lnot !A, !A \lor \lnot !A$ मिलता है। अतः
\Log{G1c} में $\Sequent !A \lor \lnot !A$ पाने के लिए !!a{proof} में समान
!!{formula} की दो उपस्थितियों को ``संकुचित'' करने की क्षमता चाहिए:
\[
\Axiom$!A \fCenter !A$
\RightLabel{\RightR{\lnot}}
\UnaryInf$\fCenter !A, \lnot !A$
\RightLabel{\RightR{\lor}}
\UnaryInf$\fCenter !A, !A \lor \lnot !A$
\RightLabel{\RightR{\lor}}
\UnaryInf$\fCenter !A \lor \lnot !A, !A \lor \lnot !A$
\RightLabel{\RightR{\Contraction}}
\UnaryInf$\fCenter !A \lor \lnot !A$
\DisplayProof
\]
वास्तव में ऐसे वैध अनुक्रम हैं जिन्हें आप दुर्बलीकरण या संकुचन के बिना
\Log{G1c} में !!{prove} नहीं कर सकते। उदाहरणतः \Log{G1c} में अनुक्रम
$!A \Sequent !B \lif !A$ का !!a{proof}, \RightR{\lif} पर समाप्त होना चाहिए
(क्योंकि केवल $!B \lif !A$ ही प्रमुख हो सकता है), और इसके लिए आधार-वाक्य
$!B, !A \Sequent !A$ चाहिए। इसे अभिगृहीति $!A \Sequent !A$ से पाने के लिए
\LeftR{\Weakening} चाहिए:
\[
\Axiom$!A \fCenter !A$
\RightLabel{\LeftR{\Weakening}}
\UnaryInf$!B, !A \fCenter !A$
\RightLabel{\RightR{\lif}}
\UnaryInf$!A \fCenter !B \lif !A$
\DisplayProof
\]

\Log{G3c} में संकुचन और दुर्बलीकरण नियमों की बिल्कुल आवश्यकता नहीं है।
दुर्बलीकरण अभिगृहीतियों में ही अंतर्निहित है। दो आधार-वाक्यों वाले नियमों में
प्रसंग $\Gamma$, $\Delta$ की दोनों प्रतियों को एक में मिलाकर संकुचन से
अधिकांशतः बचा जा सकता है। \Log{G1c} और \Log{G3c} दोनों में ऐसे ``प्रसंग
साझा करने वाले'' नियम हैं। एक आधार-वाक्य वाले नियमों में संकुचन केवल
\RightR{\lor}, \LeftR{\land}, \RightR{\lexists} और~\LeftR{\lforall} के लिए
चाहिए। \Log{G3c} में इन नियमों के रूप संकुचन को पूर्णतः अनावश्यक बना देते
हैं। एक प्रणाली \Log{G2c} (\olref{tab:G2c} देखें) है, जिसमें दो आधार-वाक्यों
वाले नियमों के प्रसंग स्वतंत्र हैं। \Log{G2c} में संकुचन \Log{G1c} से भी
अधिक महत्त्वपूर्ण है।

\subfile{rules-G2c}

\end{document}
